% ===========================================================================
%  Jàng Afrig — CM2 : CSS — Sélecteurs, Cascade, Box Model, Flexbox
%  Licence 1 DUT1-Informatique — ESP/UCAD — 2025–2026
%  Dr. El Hadji Bassirou TOURÉ
%  Compilé avec : xelatex (2–3 passes)
% ===========================================================================
\documentclass[11pt,a4paper]{article}

\usepackage{fontspec}
\setmainfont{Latin Modern Roman}
\setsansfont{Latin Modern Sans}
\setmonofont{Latin Modern Mono}
\usepackage{polyglossia}
\setdefaultlanguage{french}

\usepackage[margin=2.5cm]{geometry}
\usepackage{amsmath,amssymb}
\usepackage{graphicx}
\graphicspath{{figures/}}
\usepackage{xcolor}
\usepackage{booktabs}
\usepackage{hyperref}
\usepackage{tcolorbox}
\tcbuselibrary{theorems,breakable,skins}
\usepackage{enumitem}
\usepackage{fancyhdr}
\usepackage{array}
\usepackage{pifont}
\usepackage{longtable}
\usepackage{multirow}

% --- TikZ pour les schémas ---
\usepackage{tikz}
\usetikzlibrary{shapes.geometric,arrows.meta,positioning,calc,backgrounds,fit,shadows.blur,patterns}

\newcommand{\cmark}{\ding{51}}
\newcommand{\xmark}{\ding{55}}

\hypersetup{
    colorlinks=true,
    linkcolor=blue!60!black,
    urlcolor=blue!60!black,
    citecolor=blue!60!black
}

\newtcolorbox{objectifbox}[1][]{colback=blue!5,colframe=blue!50,fonttitle=\bfseries,title=#1,breakable}
\newtcolorbox{infobox}[1][]{colback=green!5,colframe=green!50!black,fonttitle=\bfseries,title=#1,breakable}
\newtcolorbox{attentionbox}[1][]{colback=orange!5,colframe=orange!70,fonttitle=\bfseries,title=#1,breakable}
\newtcolorbox{retenirbox}[1][]{colback=yellow!10,colframe=yellow!50!black,fonttitle=\bfseries,title=#1,breakable}
\newtcolorbox{prereqbox}[1][]{colback=cyan!5,colframe=cyan!50!black,fonttitle=\bfseries,title=#1,breakable}
\newtcolorbox{projetbox}[1][]{colback=teal!5,colframe=teal!50!black,fonttitle=\bfseries,title=#1,breakable}

% --- Couleurs box model ---
\definecolor{bmcontent}{RGB}{120,180,220}
\definecolor{bmpadding}{RGB}{180,220,140}
\definecolor{bmborder}{RGB}{220,180,90}
\definecolor{bmmargin}{RGB}{220,140,140}

% --- Couleurs Flexbox ---
\definecolor{flexitem}{RGB}{100,160,220}
\definecolor{flexaxe}{RGB}{220,80,80}
\definecolor{flexcross}{RGB}{50,140,80}

% --- Couleurs cascade ---
\definecolor{cascade1}{RGB}{220,140,140}
\definecolor{cascade2}{RGB}{220,200,140}
\definecolor{cascade3}{RGB}{140,220,180}

% --- Code source CSS et HTML ---
\usepackage{listings}
\lstdefinelanguage{HTML5}{
  morekeywords={<!DOCTYPE,html,head,meta,title,body,h1,h2,h3,h4,p,ul,ol,li,table,tr,th,td,a,img,strong,em,br,div,span,header,nav,main,section,article,footer,aside,form,input,select,option,textarea,label,button,fieldset,legend,style,link},
  sensitive=true,
  morecomment=[s]{<!--}{-->},
  morestring=[b]",
}
\lstdefinelanguage{CSS}{
  morekeywords={color,background-color,background,font-family,font-size,font-weight,text-align,padding,margin,border,border-radius,width,height,max-width,min-width,display,flex,flex-direction,justify-content,align-items,flex-wrap,gap,box-sizing,transition,transform,position,top,bottom,left,right,opacity,cursor,line-height,letter-spacing,box-shadow},
  sensitive=true,
  morecomment=[s]{/*}{*/},
  morestring=[b]",
}
\lstset{
  basicstyle=\ttfamily\small,
  backgroundcolor=\color{gray!8},
  frame=single,
  framerule=0.5pt,
  rulecolor=\color{gray!40},
  breaklines=true,
  breakatwhitespace=true,
  showstringspaces=false,
  tabsize=2,
  xleftmargin=4pt,
  xrightmargin=4pt,
  aboveskip=8pt,
  belowskip=8pt,
  literate={é}{{\'e}}1 {è}{{\`e}}1 {ê}{{\^e}}1 {à}{{\`a}}1 {ù}{{\`u}}1
           {ô}{{\^o}}1 {î}{{\^i}}1 {ç}{{\c{c}}}1 {É}{{\'E}}1
           {→}{{$\rightarrow$}}1
}

\pagestyle{fancy}
\fancyhf{}
\fancyhead[L]{\small Jàng Afrig --- CM2 : CSS}
\fancyhead[R]{\small UCAD --- L1 DUT1-Info}
\fancyfoot[C]{\thepage}
\fancyfoot[L]{\small Dr. El Hadji Bassirou TOURÉ}
\fancyfoot[R]{\small 2025--2026}
\renewcommand{\headrulewidth}{0.4pt}
\renewcommand{\footrulewidth}{0.4pt}

\title{%
\vspace{-1cm}
\textbf{\LARGE Fondamentaux du Web}\\
\textbf{\LARGE HTML, CSS, JavaScript}\\[0.8cm]
\textbf{\huge CM2 --- CSS}\\[0.3cm]
\textbf{\Large Sélecteurs, cascade, box model, Flexbox}\\[0.5cm]
\large Licence 1 DUT1-Informatique --- 2025--2026}
\author{Dr. El Hadji Bassirou TOURÉ\\
Département Génie Informatique\\
Université Cheikh Anta Diop de Dakar}
\date{}

\begin{document}

\maketitle
\thispagestyle{empty}

\vspace{0.5cm}

% ===========================================================================
% BOÎTE PRÉREQUIS
% ===========================================================================
\begin{prereqbox}[Informations pratiques]
\begin{center}
\begin{tabular}{ll}
\toprule
\textbf{Type} & Cours Magistral (CM) \\
\textbf{Séance} & S4 --- ouverture du bloc CSS \\
\textbf{Prérequis} & Lab 2 terminé (5 pages HTML poussées avec tag \texttt{v0.2}) \\
\textbf{Ce qu'on verra après} & Lab 3 : styliser Jàng Afrig (S5) \\
\textbf{Devoir} & Aucun. Lire le PDF du Lab 3 pour la prochaine séance. \\
\bottomrule
\end{tabular}
\end{center}
\end{prereqbox}

\begin{objectifbox}[Objectifs du CM]
À la fin de cette séance, vous saurez :
\begin{enumerate}[nosep]
    \item Expliquer ce que CSS \textbf{ajoute} à HTML
    \item Inclure du CSS dans une page (\texttt{<link>}, \texttt{<style>}, inline)
    \item Écrire des \textbf{sélecteurs} : élément, classe, id, descendant, pseudo-classes
    \item Comprendre la \textbf{cascade} et la \textbf{spécificité} (qui gagne quand plusieurs règles s'appliquent)
    \item Maîtriser le \textbf{Box Model} : content, padding, border, margin
    \item Utiliser \textbf{Flexbox} pour aligner et positionner des éléments
\end{enumerate}
\end{objectifbox}

\tableofcontents

\newpage

% ===========================================================================
\section{Pourquoi CSS ?}
% ===========================================================================

\subsection{Le constat du Lab 2}

À la fin du Lab 2, votre application Jàng Afrig est \textbf{complète
structurellement} : 5 pages, sémantique, formulaire, navigation. Mais
visuellement, c'est\ldots

\begin{figure}[!ht]
\centering
\begin{tikzpicture}[
    page/.style={rounded corners=4pt, draw, thick, minimum width=5.5cm, minimum height=7cm, align=left, font=\scriptsize, anchor=north west, inner sep=8pt}
]
    % Page sans CSS (gauche)
    \node[page, fill=white, draw=gray!50, label={[font=\small\bfseries, text=red!60!black]above:Sans CSS}] (sans) at (0,0) {%
        \textbf{\large Jàng Afrig}\\
        Apprendre l'Afrique, un pays à la fois\\[4pt]
        \textcolor{blue!60}{\underline{Accueil}} | \textcolor{blue!60}{\underline{Jouer}} | \textcolor{blue!60}{\underline{Mes Scores}}\\[4pt]
        \textbf{Bienvenue}\\
        Bienvenue sur \textbf{Jàng Afrig}, un quiz géographique sur les 54 pays africains.\\[4pt]
        \textbf{Les modes de jeu}\\
        \textbullet{} Capitales\\
        \textbullet{} Drapeaux\\
        \textbullet{} Langues\\
        \textbullet{} Voisins\\
        \textbullet{} Chrono\\[4pt]
        \textcolor{blue!60}{\underline{Voir le classement complet}}\\[4pt]
        © 2025--2026 Jàng Afrig
    };

    % Page avec CSS (droite)
    \node[page, fill=teal!10, draw=teal!60, line width=1pt, label={[font=\small\bfseries, text=teal!60!black]above:Avec CSS}] (avec) at (8,0) {%
        \fcolorbox{teal!50!black}{teal!30}{\parbox{4.7cm}{\centering\textbf{\large\color{white} Jàng Afrig}\\\color{white}\itshape Apprendre l'Afrique\ldots}}\\[6pt]
        \fcolorbox{gray!40}{gray!10}{\parbox{4.7cm}{\centering\small\textbf{Accueil} \textbullet{} Jouer \textbullet{} Scores}}\\[6pt]
        \textbf{\color{teal!50!black} Bienvenue}\\
        Bienvenue sur \textbf{Jàng Afrig}, un quiz géographique sur les 54 pays africains.\\[6pt]
        \textbf{\color{teal!50!black} Les modes de jeu}\\[2pt]
        \fcolorbox{orange!50!black}{orange!20}{\parbox{2.1cm}{\centering Capitales}}\hspace{2pt}
        \fcolorbox{orange!50!black}{orange!20}{\parbox{2.1cm}{\centering Drapeaux}}\\[3pt]
        \fcolorbox{orange!50!black}{orange!20}{\parbox{2.1cm}{\centering Langues}}\hspace{2pt}
        \fcolorbox{orange!50!black}{orange!20}{\parbox{2.1cm}{\centering Voisins}}\\[3pt]
        \fcolorbox{orange!50!black}{orange!20}{\parbox{2.1cm}{\centering Chrono}}\\[6pt]
        \scriptsize\itshape\color{gray} © 2025--2026 Jàng Afrig
    };
\end{tikzpicture}
\caption{La même page HTML, sans puis avec CSS.}
\label{fig:avant-apres}
\end{figure}

\textbf{Le HTML est strictement le même} dans les deux cas. Seul un fichier
CSS a été ajouté à droite. CSS, c'est la technologie qui transforme un
document brut en \textbf{interface visuelle}.

\subsection{Que peut-on faire avec CSS ?}

\begin{itemize}[leftmargin=*,nosep]
    \item Choisir des \textbf{couleurs} (texte, fond, bordures)
    \item Choisir des \textbf{polices} (Google Fonts, etc.)
    \item Définir des \textbf{tailles} et des \textbf{espacements}
    \item Positionner les éléments (alignement, grilles, colonnes)
    \item Faire du \textbf{responsive} (s'adapter au téléphone)
    \item Ajouter des \textbf{animations} et des \textbf{effets} au survol
\end{itemize}

\begin{retenirbox}[CSS = Cascading Style Sheets]
CSS signifie \textit{Cascading Style Sheets} (\og{}feuilles de style en
cascade\fg{}). \textbf{Cascading} parce que les règles CSS \emph{s'empilent}
et se combinent --- on verra ce que ça veut dire à la section sur la
cascade.
\end{retenirbox}

\newpage

% ===========================================================================
\section{Trois façons d'inclure CSS}
% ===========================================================================

CSS peut être placé à \textbf{trois endroits} différents dans une page
HTML. Chacune a ses usages et ses limites.

\subsection{1. Inline (déconseillé)}

Directement dans une balise HTML, via l'attribut \texttt{style} :

\begin{lstlisting}[language=HTML5]
<h1 style="color: red; font-size: 32px;">Jàng Afrig</h1>
\end{lstlisting}

\textbf{Quand l'utiliser :} jamais en pratique. C'est OK pour un test
rapide.

\textbf{Pourquoi déconseillé :} on mélange structure et style. Si on
change d'avis sur la couleur, il faut aller modifier \emph{chaque} balise.

\subsection{2. Interne via \texttt{<style>}}

Dans le \texttt{<head>} de la page :

\begin{lstlisting}[language=HTML5]
<head>
    <style>
        h1 { color: red; font-size: 32px; }
    </style>
</head>
\end{lstlisting}

\textbf{Quand l'utiliser :} prototype rapide, démo en classe, page unique.

\textbf{Limite :} le style ne s'applique qu'à \emph{cette} page. Pour
appliquer le même style à 5 pages, il faut copier-coller 5 fois.

\subsection{3. Externe via \texttt{<link>} (recommandé)}

Dans un fichier séparé (par convention \texttt{style.css}), lié depuis le
\texttt{<head>} :

\begin{lstlisting}[language=HTML5]
<head>
    <link rel="stylesheet" href="style.css">
</head>
\end{lstlisting}

Et dans \texttt{style.css} :

\begin{lstlisting}[language=CSS]
h1 {
    color: red;
    font-size: 32px;
}
\end{lstlisting}

\textbf{Quand l'utiliser :} \textbf{toujours, en production}.

\begin{retenirbox}[Pourquoi externe ?]
\begin{itemize}[nosep]
    \item \textbf{Réutilisable} : un seul \texttt{style.css} pour les 5 pages
    \item \textbf{Maintenable} : changer une couleur = modifier un seul endroit
    \item \textbf{Performant} : le navigateur télécharge le CSS une fois et le réutilise
    \item \textbf{Lisible} : HTML et CSS séparés, chacun son fichier
\end{itemize}
C'est la convention que vous adopterez au Lab 3.
\end{retenirbox}

\newpage

% ===========================================================================
\section{Anatomie d'une règle CSS}
% ===========================================================================

Une règle CSS, c'est toujours \textbf{la même structure} :

\begin{figure}[!ht]
\centering
\begin{tikzpicture}[
    every node/.style={font=\ttfamily\large},
    annotation/.style={font=\sffamily\small\itshape, color=red!70!black}
]
    % La règle elle-même (texte avec colorisation manuelle)
    \node (regle) {\color{blue!70!black}h1\color{black} \{ \color{purple}color\color{black}: \color{green!50!black}red\color{black}; \}};

    % Annotations
    \draw[->, color=blue!70!black, thick] (-3.0, -0.8) -- (-2.7, -0.2);
    \node[annotation, color=blue!70!black, anchor=north] at (-3.0, -0.8) {sélecteur};

    \draw[->, color=purple, thick] (-0.4, -0.8) -- (-0.4, -0.2);
    \node[annotation, color=purple, anchor=north] at (-0.4, -0.8) {propriété};

    \draw[->, color=green!50!black, thick] (1.4, -0.8) -- (1.4, -0.2);
    \node[annotation, color=green!50!black, anchor=north] at (1.4, -0.8) {valeur};

    \draw[->, color=gray, thick] (-1.5, 0.8) -- (-1.5, 0.3);
    \node[font=\sffamily\small\itshape, color=gray, anchor=south] at (-1.5, 0.8) {accolade ouvrante};

    \draw[->, color=gray, thick] (0.5, 1.4) -- (1.95, 0.3);
    \node[font=\sffamily\small\itshape, color=gray, anchor=south] at (0.5, 1.4) {fin de propriété};

    \draw[->, color=gray, thick] (2.6, 0.8) -- (2.6, 0.3);
    \node[font=\sffamily\small\itshape, color=gray, anchor=south] at (2.6, 0.8) {accolade fermante};
\end{tikzpicture}
\caption{Anatomie d'une règle CSS.}
\label{fig:anatomie-css}
\end{figure}

\subsection{Une règle plus complète}

\begin{lstlisting}[language=CSS]
h1 {
    color: red;
    font-size: 32px;
    text-align: center;
    margin-bottom: 20px;
}
\end{lstlisting}

\begin{itemize}[nosep]
    \item \textbf{Sélecteur} : \texttt{h1} --- cible \emph{tous les \texttt{<h1>}} de la page
    \item \textbf{Bloc} : entre \texttt{\{} et \texttt{\}}
    \item \textbf{Propriétés} : \texttt{color}, \texttt{font-size}, \texttt{text-align}, \texttt{margin-bottom}
    \item Chaque \textbf{déclaration} (propriété : valeur) finit par un \texttt{;}
\end{itemize}

\begin{retenirbox}[Le point-virgule]
Le \texttt{;} après chaque déclaration est \textbf{techniquement} optionnel
pour la dernière, mais \textbf{toujours mettez-le}. Vous oublierez d'en
mettre un quand vous ajouterez une nouvelle propriété, et vous chercherez
le bug pendant 20 minutes.
\end{retenirbox}

\newpage

% ===========================================================================
\section{Les sélecteurs}
% ===========================================================================

Les \textbf{sélecteurs} disent \og{}à quels éléments cette règle
s'applique\fg{}. CSS en propose plusieurs dizaines, mais 5 vous
suffiront pour 90\% des cas.

\subsection{Sélecteur d'élément}

Cible \textbf{toutes les balises de ce type} :

\begin{lstlisting}[language=CSS]
p {
    color: gray;
}
\end{lstlisting}

$\to$ tous les \texttt{<p>} de toutes les pages liant ce CSS deviennent
gris.

\subsection{Sélecteur de classe (le plus utilisé)}

Cible toutes les balises ayant un attribut \texttt{class="..."}. On
préfixe le nom par un point :

HTML :
\begin{lstlisting}[language=HTML5]
<button class="btn-primary">Commencer</button>
<button class="btn-secondary">Annuler</button>
\end{lstlisting}

CSS :
\begin{lstlisting}[language=CSS]
.btn-primary {
    background-color: green;
    color: white;
}

.btn-secondary {
    background-color: gray;
    color: white;
}
\end{lstlisting}

\textbf{Une classe peut être utilisée sur plusieurs éléments}. C'est ça
qui fait sa puissance.

\subsection{Sélecteur d'identifiant}

Cible la balise ayant un attribut \texttt{id="..."}. Préfixé par
\texttt{\#} :

HTML :
\begin{lstlisting}[language=HTML5]
<header id="entete-principal">
    <h1>Jàng Afrig</h1>
</header>
\end{lstlisting}

CSS :
\begin{lstlisting}[language=CSS]
#entete-principal {
    background-color: teal;
}
\end{lstlisting}

\begin{attentionbox}[\texttt{id} vs. \texttt{class}]
\begin{itemize}[nosep]
    \item Un \texttt{id} doit être \textbf{unique} sur la page (un seul \texttt{\#entete-principal})
    \item Une \texttt{class} peut être \textbf{réutilisée} (plusieurs \texttt{.btn-primary})
\end{itemize}
\textbf{En pratique, on utilise presque toujours des classes.} Les \texttt{id}
servent surtout à JS plus tard.
\end{attentionbox}

\subsection{Sélecteur descendant}

Cible une balise \emph{à l'intérieur d'une autre} :

\begin{lstlisting}[language=CSS]
nav a {
    color: white;
    text-decoration: none;
}
\end{lstlisting}

$\to$ tous les \texttt{<a>} qui sont \emph{à l'intérieur} d'un \texttt{<nav>}
deviennent blancs et sans soulignement. Les autres \texttt{<a>} de la
page ne sont pas touchés.

\subsection{Pseudo-classes}

Cible un élément \emph{dans un certain état} :

\begin{lstlisting}[language=CSS]
button:hover {
    background-color: orange;
}

a:visited {
    color: purple;
}

input:focus {
    border-color: green;
}
\end{lstlisting}

\begin{itemize}[nosep]
    \item \texttt{:hover} : quand la souris passe dessus
    \item \texttt{:focus} : quand l'élément est sélectionné (au clavier ou cliqué)
    \item \texttt{:visited} : pour un lien déjà cliqué
\end{itemize}

\begin{infobox}[Démo en direct]
Le prof va montrer en direct, dans la console DevTools, comment ces
sélecteurs s'appliquent à la page d'accueil de Jàng Afrig.
\end{infobox}

\newpage

% ===========================================================================
\section{Cascade et spécificité}
% ===========================================================================

\subsection{Le problème}

Que se passe-t-il quand \textbf{deux règles différentes} essaient de styler
\textbf{le même élément} avec \textbf{deux valeurs incompatibles} ?

\begin{lstlisting}[language=CSS]
p { color: red; }
p { color: blue; }
\end{lstlisting}

\textbf{Question :} les paragraphes seront rouges ou bleus ?

\subsection{Règle 1 : la cascade (le dernier gagne)}

Si deux règles ont la \textbf{même spécificité}, c'est \textbf{la dernière
écrite} qui gagne.

\begin{figure}[!ht]
\centering
\begin{tikzpicture}[
    couche/.style={rounded corners=4pt, draw, thick, minimum width=8cm, minimum height=1cm, align=center, font=\small}
]
    \node[couche, fill=cascade1!40, draw=cascade1!70!black] (c1) {\texttt{p \{ color: red; \}}};
    \node[couche, fill=cascade2!40, draw=cascade2!70!black, below=0.1cm of c1] (c2) {\texttt{p \{ color: blue; \}} \quad\textit{(plus bas $\to$ gagne)}};
    \node[couche, fill=cascade3!40, draw=cascade3!70!black, below=0.6cm of c2] (resultat) {Résultat affiché : \textbf{\color{blue!70!black}bleu}};

    \draw[-{Stealth[length=3mm]}, thick, color=gray] (c2.south) -- (resultat.north);
\end{tikzpicture}
\caption{Quand deux règles de même spécificité s'opposent, la dernière écrite gagne.}
\label{fig:cascade}
\end{figure}

\subsection{Règle 2 : la spécificité (le plus précis gagne)}

Mais si une règle est \textbf{plus précise}, elle l'emporte, peu importe
l'ordre. La spécificité se calcule comme suit :

\begin{center}
\begin{tabular}{@{}lcl@{}}
\toprule
\textbf{Type de sélecteur} & \textbf{Score} & \textbf{Exemple} \\
\midrule
Sélecteur d'élément       & 1   & \texttt{p}              \\
Sélecteur descendant      & 2 (1+1) & \texttt{nav a}     \\
Sélecteur de classe       & 10  & \texttt{.btn-primary}   \\
Sélecteur d'identifiant   & 100 & \texttt{\#entete-principal} \\
\texttt{style="..."} inline & 1000 & --- \\
\bottomrule
\end{tabular}
\end{center}

\subsection{Exemple chiffré}

\begin{lstlisting}[language=CSS]
p { color: red; }                      /* score 1   */
.intro { color: blue; }                /* score 10  */
#message-special { color: green; }     /* score 100 */
\end{lstlisting}

Pour ce HTML :
\begin{lstlisting}[language=HTML5]
<p class="intro" id="message-special">Texte</p>
\end{lstlisting}

Le texte sera\ldots \textbf{vert} (le \texttt{\#message-special} a le score
le plus élevé).

\begin{retenirbox}[Conseil pratique]
\begin{itemize}[nosep]
    \item \textbf{Préférez les classes} (\texttt{.}) aux ids (\texttt{\#})
    \item \textbf{Évitez les sélecteurs trop précis} comme \texttt{header nav ul li a}
    \item \textbf{N'utilisez jamais \texttt{style="..."} inline} sauf en dernier recours
\end{itemize}
Si vous gardez des règles simples (1-2 niveaux de précision), vous éviterez
99\% des bugs de spécificité.
\end{retenirbox}

\newpage

% ===========================================================================
\section{Le Box Model}
% ===========================================================================

\textbf{Le Box Model est sans doute le concept CSS le plus important.}
Tout élément HTML est une \textbf{boîte}, et chaque boîte a 4 zones.

\subsection{Les 4 zones}

\begin{figure}[!ht]
\centering
\begin{tikzpicture}[
    every node/.style={font=\sffamily\bfseries}
]
    % Margin
    \filldraw[fill=bmmargin!40, draw=bmmargin!70!black, thick, rounded corners=2pt] (0,0) rectangle (10,7);
    \node[anchor=north west, color=bmmargin!50!black, font=\sffamily\bfseries\small] at (0.1, 6.9) {margin};

    % Border
    \filldraw[fill=bmborder!50, draw=bmborder!70!black, thick] (1,1) rectangle (9,6);
    \node[anchor=north west, color=bmborder!50!black, font=\sffamily\bfseries\small] at (1.1, 5.9) {border};

    % Padding
    \filldraw[fill=bmpadding!40, draw=bmpadding!70!black, thick] (1.5,1.5) rectangle (8.5,5.5);
    \node[anchor=north west, color=bmpadding!50!black, font=\sffamily\bfseries\small] at (1.6, 5.4) {padding};

    % Content
    \filldraw[fill=bmcontent!50, draw=bmcontent!70!black, thick] (3,2.5) rectangle (7,4.5);
    \node[font=\sffamily\bfseries, color=white, text=black] at (5, 3.5) {content};

    % Légendes externes
    \node[font=\sffamily\small\itshape, color=gray] at (5, -0.4) {(la zone visible totale)};
\end{tikzpicture}
\caption{Le Box Model : content au centre, entouré de padding, border, margin.}
\label{fig:boxmodel}
\end{figure}

\subsection{Ce que fait chaque zone}

\begin{center}
\begin{tabular}{@{}lp{8cm}@{}}
\toprule
\textbf{Zone}    & \textbf{Rôle} \\
\midrule
\textbf{content} & Le contenu réel : texte, image, autre balise. Sa taille est définie par \texttt{width} et \texttt{height}. \\
\addlinespace
\textbf{padding} & Espace \emph{intérieur} entre le contenu et la bordure. \og{}Du rembourrage à l'intérieur de la boîte\fg{}. \\
\addlinespace
\textbf{border}  & La bordure visible (épaisseur, style, couleur). \\
\addlinespace
\textbf{margin}  & Espace \emph{extérieur} entre cette boîte et les boîtes voisines. \og{}Du vide autour de la boîte\fg{}. \\
\bottomrule
\end{tabular}
\end{center}

\begin{infobox}[L'analogie du paquet cadeau]
\begin{itemize}[nosep]
    \item \textbf{content} : le cadeau lui-même
    \item \textbf{padding} : le papier de soie autour, dans la boîte
    \item \textbf{border} : le carton de la boîte
    \item \textbf{margin} : l'espace que vous laissez \emph{autour} pour ne pas écraser les autres paquets
\end{itemize}
\end{infobox}

\subsection{Exemple en CSS}

\begin{lstlisting}[language=CSS]
.carte-mode {
    width: 200px;
    padding: 20px;
    border: 2px solid teal;
    margin: 10px;
    background-color: #f4f4f4;
}
\end{lstlisting}

\subsection{Le piège classique}

Par défaut, \texttt{width} ne mesure que le \textbf{content}. Donc une boîte
de 200px avec \texttt{padding: 20px} occupe en réalité \textbf{240px} de
large (200 + 20 + 20).

\textbf{Solution} : ajouter cette règle au début de votre CSS pour que
\texttt{width} inclue padding et border :

\begin{lstlisting}[language=CSS]
* {
    box-sizing: border-box;
}
\end{lstlisting}

\begin{retenirbox}[Réflexe pro]
\textbf{Toujours mettre \texttt{box-sizing: border-box}} en début de
fichier CSS. C'est le réflexe de tous les développeurs web modernes. Sans
ça, vos calculs de largeur seront faux et vous galérerez en responsive.
\end{retenirbox}

\subsection{Démo en direct dans DevTools}

\begin{infobox}[Le Box Model dans DevTools]
Le prof va montrer en direct comment voir le Box Model d'un élément :
\begin{enumerate}[nosep]
    \item Clic droit sur l'élément $\to$ \textbf{Inspecter}
    \item Onglet \textbf{Computed} ou panneau de droite
    \item Le Box Model s'affiche avec les valeurs réelles de chaque zone
\end{enumerate}
C'est l'outil que vous utiliserez \emph{tout le temps} pour débugger
votre CSS.
\end{infobox}

\newpage

% ===========================================================================
\section{Flexbox : aligner et positionner}
% ===========================================================================

Avant Flexbox, centrer un bouton verticalement dans une page demandait
des \emph{hacks}. Avec Flexbox, c'est trois lignes.

\subsection{Le principe}

Flexbox transforme un élément parent en \textbf{conteneur flexible} (un
\og{}flex container\fg{}). Ses enfants directs deviennent des
\textbf{flex items} qui peuvent être alignés et distribués librement.

\begin{lstlisting}[language=CSS]
.conteneur {
    display: flex;
}
\end{lstlisting}

C'est tout. Le parent passe en mode flex, et ses enfants se mettent
\textbf{en ligne horizontale} par défaut.

\subsection{Les deux axes}

\begin{figure}[!ht]
\centering
\begin{tikzpicture}[
    item/.style={rounded corners=3pt, fill=flexitem!40, draw=flexitem!70!black, thick, minimum width=1.5cm, minimum height=1cm, font=\small\bfseries},
    axe/.style={->, thick, line width=1.2pt}
]
    % Conteneur flex
    \draw[draw=gray!50, thick, dashed, fill=gray!5, rounded corners=4pt] (-0.3, -0.7) rectangle (8.3, 1.7);

    % Items
    \node[item] (i1) at (0.85, 0.5) {1};
    \node[item] (i2) at (2.65, 0.5) {2};
    \node[item] (i3) at (4.45, 0.5) {3};
    \node[item] (i4) at (6.25, 0.5) {4};

    % Axe principal (horizontal)
    \draw[axe, color=flexaxe] (-0.3, -1.3) -- (8.5, -1.3);
    \node[font=\small\bfseries, color=flexaxe, anchor=west] at (-0.3, -1.7) {axe principal};
    \node[font=\scriptsize\itshape, color=flexaxe, anchor=west] at (-0.3, -2.1) {(piloté par \texttt{justify-content})};

    % Axe secondaire (vertical)
    \draw[axe, color=flexcross] (10, -0.7) -- (10, 1.7);
    \node[font=\small\bfseries, color=flexcross, anchor=west] at (10.2, 1.4) {axe secondaire};
    \node[font=\scriptsize\itshape, color=flexcross, anchor=west] at (10.2, 0.9) {(piloté par};
    \node[font=\scriptsize\itshape, color=flexcross, anchor=west] at (10.2, 0.5) {\texttt{align-items})};
\end{tikzpicture}
\caption{Les deux axes de Flexbox : principal (horizontal) et secondaire (vertical).}
\label{fig:flex-axes}
\end{figure}

\subsection{Les propriétés essentielles}

\subsubsection*{Sur le conteneur :}

\begin{center}
\begin{tabular}{@{}lp{8cm}@{}}
\toprule
\textbf{Propriété}        & \textbf{Rôle} \\
\midrule
\texttt{display: flex}    & Active Flexbox sur le conteneur \\
\texttt{flex-direction}   & \texttt{row} (par défaut, horizontal) ou \texttt{column} (vertical) \\
\texttt{justify-content}  & Aligne les items sur \textbf{l'axe principal} \\
\texttt{align-items}      & Aligne les items sur \textbf{l'axe secondaire} \\
\texttt{flex-wrap}        & \texttt{wrap} : retour à la ligne si pas la place \\
\texttt{gap}              & Espace entre les items \\
\bottomrule
\end{tabular}
\end{center}

\subsection{Les valeurs de \texttt{justify-content}}

\begin{center}
\begin{tabular}{@{}ll@{}}
\toprule
\textbf{Valeur}      & \textbf{Effet} \\
\midrule
\texttt{flex-start}    & Items collés à gauche (par défaut) \\
\texttt{center}        & Items centrés horizontalement \\
\texttt{flex-end}      & Items collés à droite \\
\texttt{space-between} & Premier à gauche, dernier à droite, le reste réparti \\
\texttt{space-around}  & Espacement égal autour de chaque item \\
\bottomrule
\end{tabular}
\end{center}

\subsection{Centrer un bouton (le grand classique)}

\begin{lstlisting}[language=CSS]
.parent {
    display: flex;
    justify-content: center;   /* centré horizontalement */
    align-items: center;       /* centré verticalement */
    height: 200px;
}
\end{lstlisting}

3 lignes pour centrer parfaitement, à droite ET en bas, à hauteur ET en
largeur. Avant Flexbox (avant 2015), c'était un cauchemar.

\subsection{Cas d'usage typique : la nav}

\begin{lstlisting}[language=HTML5]
<nav class="nav-principale">
    <a href="index.html">Accueil</a>
    <a href="quiz.html">Jouer</a>
    <a href="scores.html">Mes Scores</a>
    <a href="apropos.html">À propos</a>
</nav>
\end{lstlisting}

\begin{lstlisting}[language=CSS]
.nav-principale {
    display: flex;
    justify-content: space-between;
    padding: 10px 20px;
    background-color: teal;
}

.nav-principale a {
    color: white;
    text-decoration: none;
}
\end{lstlisting}

$\to$ une barre de navigation propre avec liens espacés et fond coloré, en
quelques lignes.

\begin{retenirbox}[Mémo Flexbox]
Les 5 propriétés à connaître :
\begin{enumerate}[nosep]
    \item \texttt{display: flex} \quad (active Flexbox)
    \item \texttt{justify-content} \quad (axe principal)
    \item \texttt{align-items} \quad (axe secondaire)
    \item \texttt{flex-wrap: wrap} \quad (retour à la ligne)
    \item \texttt{gap} \quad (espacement entre items)
\end{enumerate}
Avec ces 5, vous gérez 95\% des layouts d'application.
\end{retenirbox}

\newpage

% ===========================================================================
\section{Mini-exercice en classe}
% ===========================================================================

\begin{objectifbox}[Exercice guidé --- on le fait ensemble]

\textbf{Énoncé.} Sur la page d'accueil de Jàng Afrig, ajoutons une règle
CSS qui :
\begin{enumerate}[nosep]
    \item Met le \texttt{<header>} sur fond \textbf{teal} avec texte \textbf{blanc}
    \item Aligne les liens du \texttt{<nav>} en \textbf{ligne horizontale}, espacés \textbf{equitablement}
    \item Centre le bouton «\,Commencer\,» horizontalement dans son conteneur
\end{enumerate}

\medskip

\textbf{On va l'écrire ensemble} dans un \texttt{<style>} interne pour
aller vite (au Lab 3, on basculera en fichier externe).
\end{objectifbox}

\subsection{Solution attendue}

\begin{lstlisting}[language=CSS]
/* Box-sizing pour tout */
* { box-sizing: border-box; }

/* Header coloré */
header {
    background-color: teal;
    color: white;
    padding: 20px;
    text-align: center;
}

/* Nav en flex */
nav {
    display: flex;
    justify-content: space-around;
    background-color: #006644;
    padding: 10px;
}

nav a {
    color: white;
    text-decoration: none;
}

/* Bouton centré */
.bouton-conteneur {
    display: flex;
    justify-content: center;
    padding: 20px;
}
\end{lstlisting}

\begin{infobox}[Si vous bloquez]
Pas de panique. Ce sera \emph{exactement} ce qu'on fait au Lab 3, en
détail, étape par étape. L'objectif aujourd'hui est de \textbf{voir la
mécanique}, pas de tout mémoriser.
\end{infobox}

\newpage

% ===========================================================================
\section{Résumé et transition}
% ===========================================================================

\subsection{Ce qu'on a vu}

\begin{itemize}[leftmargin=*,nosep]
    \item CSS = \textbf{Cascading Style Sheets} : la couche visuelle du Web
    \item \textbf{3 façons d'inclure} CSS : inline, \texttt{<style>}, \texttt{<link>} (recommandé)
    \item Anatomie d'une règle : \texttt{sélecteur \{ propriété: valeur; \}}
    \item \textbf{5 sélecteurs} : élément, classe (\texttt{.}), id (\texttt{\#}), descendant, pseudo-classes (\texttt{:hover})
    \item \textbf{Cascade} : à spécificité égale, le dernier gagne
    \item \textbf{Spécificité} : id (100) > classe (10) > élément (1)
    \item \textbf{Box Model} : content + padding + border + margin
    \item Réflexe : toujours \texttt{box-sizing: border-box}
    \item \textbf{Flexbox} : 5 propriétés pour aligner et positionner
\end{itemize}

\subsection{Ce qu'on fera au Lab 3}

\begin{itemize}[leftmargin=*,nosep]
    \item Créer le fichier \texttt{style.css} dans le repo
    \item Le lier aux 5 pages avec \texttt{<link>}
    \item Définir la \textbf{palette panafricaine} (vert, jaune, rouge + neutre)
    \item Utiliser une \textbf{police Google Fonts}
    \item Styliser le \texttt{<header>}, le \texttt{<nav>}, les sections, le footer
    \item Créer des \textbf{boutons} stylés avec \texttt{:hover}
    \item Tag \texttt{v0.3}
\end{itemize}

\subsection{Consigne pour la prochaine séance}

\begin{attentionbox}[Avant le Lab 3]
\begin{enumerate}[nosep]
    \item \textbf{Lisez} le PDF du Lab 3.
    \item \textbf{Vérifiez} que vos 5 pages HTML s'affichent toujours correctement (\texttt{git pull} d'abord si vous avez travaillé sur un autre ordinateur).
    \item \textbf{Choisissez 1 ou 2 polices} sur \href{https://fonts.google.com}{Google Fonts} qui vous plaisent. Idée : une «\,sans-serif\,» moderne pour le texte (\emph{Poppins}, \emph{Inter}, \emph{Nunito}\ldots) et une plus marquée pour les titres.
\end{enumerate}
Aucun code à écrire à la maison.
\end{attentionbox}

\vspace{0.5cm}

\begin{center}
\fbox{
\begin{minipage}{0.85\textwidth}
\centering
\vspace{0.5cm}
\textbf{\Large CM2 --- Terminé \cmark}\\[0.3cm]
\normalsize
Vous savez comment CSS donne vie au HTML.\\
Prochaine étape : \textbf{habiller Jàng Afrig}.\\[0.2cm]
\textit{Rendez-vous au Lab 3 !}\\[0.3cm]
\vspace{0.3cm}
\end{minipage}
}
\end{center}

\end{document}
